\section{Results and Evaluation}

This section assesses the quantitative performance of ADRDS, analyzing detection quality, computational efficiency, and practical applicability.

\subsection{Supervised Classifier Performance}

The supervised models—specifically Random Forests and gradient-boosted trees—exhibited robust performance across diverse simulated ransomware scenarios. Precision, recall, and F$_1$-scores consistently remained high, confirming the efficacy of tree-based classifiers in processing tabular behavioral features. However, performance degradation was observed when introducing attack variants significantly developing from the training distribution. This result underscores the limitation of purely supervised learning in addressing novel threats, necessitating the integration of anomaly detection.

\subsection{Anomaly Detection Results}

The LSTM-based autoencoder demonstrated significant utility in identifying behaviors undetected by supervised models. Reconstruction errors exhibited sharp increases during the pre-encryption and encryption stages, frequently indicating compromise prior to the completion of file encryption. Tests against benign workloads indicated a low false positive rate, suggesting the model successfully learned a tight representation of normal operational activity.

\subsection{Hybrid Risk Fusion}

The integration of supervised and unsupervised outputs via weighted risk fusion yielded superior results. Compared to isolated supervised detection, the fusion approach reduced false positives. Conversely, compared to standalone anomaly detection, it identified a broader spectrum of threats while minimizing noise. Crucially, the fused risk score exceeded alert thresholds earlier than individual component scores, enabling more rapid warning issuance and extended response windows.

\subsection{System Latency}

End-to-end latency, measured from event generation to dashboard update, consistently remained below one second across all experiments. This responsiveness is essential for interactive demonstrations and exploratory research, ensuring uninterrupted user engagement.

\subsection{Conclusion of Results}

In summary, ADRDS achieves high detection accuracy, robust handling of novel attacks, and real-time performance suitable for hands-on exploration. Supervised models address known threats, the autoencoder identifies anomalies, and risk fusion integrates these inputs into a actionable metric. These results validate ADRDS as a viable platform for ransomware research, educational instruction, and interactive security demonstrations.
